\chapter{الدوال الحسابية والبنية الضربية}

\section{مقدمة}

الدالة الحسابية هي دالة معرفة على الأعداد الصحيحة الموجبة وقيمها في مجموعة عددية، وغالبًا في \(\mathbb{C}\). وتظهر هذه الدوال عندما نريد ترميز خاصية حسابية لكل عدد صحيح في قيمة واحدة.

من أمثلتها:

\[
1(n)=1,
\qquad
\operatorname{id}(n)=n,
\qquad
\mu(n),
\qquad
\varphi(n),
\qquad
\Lambda(n),
\qquad
\tau(n).
\]

القوة الحقيقية لهذه الدوال لا تأتي من قيمها منفردة، بل من العلاقات التي تربطها تحت عمليات مثل الالتفاف الديريشلي، ومن قدرتها على عكس البنية الضربية للأعداد الصحيحة.

\section{فضاء الدوال الحسابية}

لنرمز إلى مجموعة جميع الدوال الحسابية ذات القيم العقدية بالرمز
\[
\mathcal{A}
=
\{f:\mathbb{N}\to\mathbb{C}\}.
\]

نعرف الجمع والضرب النقطي بواسطة
\[
(f+g)(n)=f(n)+g(n),
\]
و
\[
(fg)(n)=f(n)g(n).
\]

بهاتين العمليتين تكون \(\mathcal{A}\) جبرًا تبديليًا على \(\mathbb{C}\). لكن العملية الأهم في نظرية الأعداد التحليلية ليست الضرب النقطي، بل الالتفاف الديريشلي.

\section{الدوال الضربية}

\begin{definition}
تسمى الدالة الحسابية \(f\) ضربية إذا تحقق:

\begin{enumerate}
  \item \(f(1)=1\).
  \item لكل \(m,n\) أوليين فيما بينهما،
  \[
  f(mn)=f(m)f(n).
  \]
\end{enumerate}
\end{definition}

\begin{definition}
تسمى \(f\) ضربية تمامًا إذا
\[
f(mn)=f(m)f(n)
\]
لكل \(m,n\geq1\)، دون اشتراط \((m,n)=1\).
\end{definition}

كل دالة ضربية تمامًا ضربية، لكن العكس غير صحيح.

\begin{example}
الدالة
\[
f(n)=n^s
\]
ضربية تمامًا لكل \(s\in\mathbb{C}\).
\end{example}

\begin{example}
دالة أويلر \(\varphi\) ضربية، لكنها ليست ضربية تمامًا، لأن
\[
\varphi(4)=2
\]
بينما
\[
\varphi(2)\varphi(2)=1.
\]
\end{example}

\section{تحديد الدالة الضربية على القوى الأولية}

\begin{theorem}
إذا كانت \(f\) ضربية، فإن قيمها تتحدد تمامًا من قيمها على القوى الأولية
\[
f(p^k).
\]
\end{theorem}

\begin{proof}
إذا كان
\[
n=\prod_{j=1}^{r}p_j^{\alpha_j}
\]
هو التحليل الأولي لـ\(n\)، فإن القوى الأولية المختلفة أولية فيما بينها. وبالضربية:
\[
f(n)
=
\prod_{j=1}^{r}f(p_j^{\alpha_j}).
\]
\end{proof}

إذن دراسة الدالة الضربية تنقسم طبيعيًا إلى دراسة محلية عند كل عدد أولي.

\begin{remark}
هذا هو الأصل الحسابي لفكرة العوامل المحلية في المنتجات الأويلرية ودوال \(L\).
\end{remark}

\section{الالتفاف الديريشلي}

\begin{definition}
إذا كانت \(f,g\in\mathcal{A}\)، فإن التفافهما الديريشلي هو الدالة
\[
(f*g)(n)
=
\sum_{d\mid n}f(d)g\left(\frac{n}{d}\right).
\]
\end{definition}

ويمكن كتابة الصيغة بصورة متناظرة:
\[
(f*g)(n)
=
\sum_{ab=n}f(a)g(b).
\]

\begin{example}
إذا كانت \(1(n)=1\)، فإن
\[
(1*1)(n)
=
\sum_{d\mid n}1
=
\tau(n),
\]
حيث \(\tau(n)\) عدد قواسم \(n\).
\end{example}

\begin{example}
إذا كانت \(\operatorname{id}(n)=n\)، فإن
\[
(1*\operatorname{id})(n)
=
\sum_{d\mid n}\frac{n}{d}
=
\sum_{d\mid n}d
=
\sigma(n).
\]
\end{example}

\section{الخواص الجبرية للالتفاف}

\begin{theorem}
الالتفاف الديريشلي تبديلي وتجميعي وتوزيعي على الجمع.
\end{theorem}

\begin{proof}
التبديلية تتبع من تبديل \(d\) مع \(n/d\):
\[
(f*g)(n)
=
\sum_{d\mid n}f(d)g(n/d)
=
\sum_{d\mid n}g(d)f(n/d)
=
(g*f)(n).
\]

أما التجميعية:
\[
((f*g)*h)(n)
=
\sum_{abc=n}f(a)g(b)h(c)
=
(f*(g*h))(n).
\]

والتوزيعية تتبع مباشرة من خطية المجموع.
\end{proof}

\section{عنصر الوحدة}

\begin{definition}
نعرف الدالة \(\varepsilon\) بواسطة
\[
\varepsilon(1)=1,
\qquad
\varepsilon(n)=0
\quad(n>1).
\]
\end{definition}

\begin{proposition}
لكل \(f\in\mathcal{A}\):
\[
f*\varepsilon=\varepsilon*f=f.
\]
\end{proposition}

\begin{proof}
لدينا
\[
(f*\varepsilon)(n)
=
\sum_{d\mid n}f(d)\varepsilon(n/d).
\]
الحد الوحيد غير الصفري يحدث عندما \(n/d=1\)، أي \(d=n\). لذا
\[
(f*\varepsilon)(n)=f(n).
\]
\end{proof}

إذن \((\mathcal{A},+,*)\) حلقة تبادلية ذات وحدة.

\section{المعكوس الديريشلي}

\begin{definition}
تسمى \(g\) معكوسًا ديريشليًا لـ\(f\) إذا
\[
f*g=\varepsilon.
\]
\end{definition}

\begin{theorem}
تملك الدالة \(f\) معكوسًا ديريشليًا إذا وفقط إذا
\[
f(1)\neq0.
\]
\end{theorem}

\begin{proof}
إذا كان \(f*g=\varepsilon\)، فإن
\[
f(1)g(1)=1،
\]
ومن ثم \(f(1)\neq0\).

عكسيًا، إذا كان \(f(1)\neq0\)، نعرف \(g\) استقرائيًا. أولًا:
\[
g(1)=\frac1{f(1)}.
\]
ولكل \(n>1\)، يجب أن يكون
\[
0=(f*g)(n)
=
f(1)g(n)
+
\sum_{\substack{d\mid n\\d>1}}f(d)g(n/d).
\]
إذن
\[
g(n)
=
-\frac1{f(1)}
\sum_{\substack{d\mid n\\d>1}}f(d)g(n/d).
\]
وهذا يحدد \(g(n)\) من قيم سابقة.
\end{proof}

\section{دالة موبيوس}

\begin{definition}
تعرف دالة موبيوس \(\mu\) بواسطة
\[
\mu(n)=
\begin{cases}
1,&n=1,\\
(-1)^r,&n \text{ حاصل ضرب } r \text{ أوليات متميزة},\\
0,&p^2\mid n \text{ لبعض الأوليات }p.
\end{cases}
\]
\end{definition}

\begin{example}
\[
\mu(1)=1,\qquad
\mu(6)=1,\qquad
\mu(30)=-1,\qquad
\mu(12)=0.
\]
\end{example}

\begin{proposition}
الدالة \(\mu\) ضربية.
\end{proposition}

\begin{proof}
إذا كان \((m,n)=1\)، فإن مجموع العوامل الأولية المتميزة في \(mn\) هو مجموعها في \(m\) و\(n\). كما أن \(mn\) خال من المربعات إذا وفقط إذا كان كل من \(m,n\) خاليًا من المربعات. ومن ثم
\[
\mu(mn)=\mu(m)\mu(n).
\]
\end{proof}

\section{الهوية الأساسية لموبيوس}

\begin{theorem}
لدينا
\[
\sum_{d\mid n}\mu(d)
=
\begin{cases}
1,&n=1,\\
0,&n>1.
\end{cases}
\]
أي
\[
\mu*1=\varepsilon.
\]
\end{theorem}

\begin{proof}
إذا كان
\[
n=p_1^{\alpha_1}\cdots p_r^{\alpha_r},
\]
فإن القواسم \(d\) التي تسهم في المجموع هي فقط القواسم الخالية من المربعات، أي
\[
d=\prod_{j\in S}p_j
\]
لمجموعة جزئية \(S\subseteq\{1,\dots,r\}\). ومن ثم
\[
\sum_{d\mid n}\mu(d)
=
\sum_{S\subseteq\{1,\dots,r\}}(-1)^{|S|}
=
(1-1)^r.
\]
إذا \(n>1\)، فـ\(r\geq1\) ويكون المجموع صفرًا. أما \(n=1\)، فالمجموع يساوي \(1\).
\end{proof}

إذن \(\mu\) هي المعكوس الديريشلي للدالة الثابتة \(1\).

\section{صيغة قلب موبيوس}

\begin{theorem}[قلب موبيوس]
إذا
\[
F(n)=\sum_{d\mid n}f(d),
\]
فإن
\[
f(n)=\sum_{d\mid n}\mu(d)F(n/d).
\]
\end{theorem}

\begin{proof}
الفرضية تعني
\[
F=1*f.
\]
بالتفاف الطرفين مع \(\mu\):
\[
\mu*F
=
\mu*(1*f)
=
(\mu*1)*f
=
\varepsilon*f
=
f.
\]
\end{proof}

\begin{corollary}
إذا
\[
F(n)=\sum_{d\mid n}f(n/d),
\]
فإن
\[
f(n)=\sum_{d\mid n}\mu(d)F(n/d).
\]
\end{corollary}

\section{مثال: استرجاع دالة أويلر}

من الهوية
\[
n=\sum_{d\mid n}\varphi(d)
\]
نحصل بقلب موبيوس على
\[
\varphi(n)
=
\sum_{d\mid n}\mu(d)\frac{n}{d}.
\]

أي
\[
\varphi
=
\mu*\operatorname{id}.
\]

وبالتالي:
\[
\varphi(n)
=
n\sum_{d\mid n}\frac{\mu(d)}{d}.
\]

إذا كان
\[
n=\prod_{p\mid n}p^{\alpha_p},
\]
فإن
\[
\varphi(n)
=
n\prod_{p\mid n}\left(1-\frac1p\right).
\]

\section{حفظ الضربية تحت الالتفاف}

\begin{theorem}
إذا كانت \(f\) و\(g\) دالتين ضربيتين، فإن \(f*g\) ضربية.
\end{theorem}

\begin{proof}
لنفرض \((m,n)=1\). كل قاسم \(d\mid mn\) يكتب بصورة وحيدة:
\[
d=d_1d_2,
\qquad
d_1\mid m,
\quad
d_2\mid n.
\]
وعندئذ:
\[
\frac{mn}{d}
=
\frac{m}{d_1}\frac{n}{d_2}.
\]

وباستخدام الضربية:
\[
(f*g)(mn)
=
\sum_{d_1\mid m}
\sum_{d_2\mid n}
f(d_1d_2)
g\left(\frac{m}{d_1}\frac{n}{d_2}\right).
\]
إذن
\[
(f*g)(mn)
=
\left(\sum_{d_1\mid m}
f(d_1)g(m/d_1)\right)
\left(\sum_{d_2\mid n}
f(d_2)g(n/d_2)\right),
\]
أي
\[
(f*g)(mn)
=
(f*g)(m)(f*g)(n).
\]
\end{proof}

\begin{corollary}
إذا كانت \(f\) ضربية ولها معكوس ديريشلي، فإن \(f^{-1}\) ضربية.
\end{corollary}

\section{الالتفاف على القوى الأولية}

إذا كانت \(f,g\) ضربيتين، فإن
\[
(f*g)(p^\alpha)
=
\sum_{j=0}^{\alpha}
f(p^j)g(p^{\alpha-j}).
\]

هذه صيغة التفاف عادية للمتتاليتين المحليتين
\[
\{f(p^j)\}_{j\geq0},
\qquad
\{g(p^j)\}_{j\geq0}.
\]

ومن هنا تظهر متسلسلات بيل المحلية:
\[
\mathcal{B}_p(f;X)
=
\sum_{j=0}^{\infty}f(p^j)X^j.
\]

وعندئذ:
\[
\mathcal{B}_p(f*g;X)
=
\mathcal{B}_p(f;X)\mathcal{B}_p(g;X).
\]

\section{الضرب النقطي والالتفاف}

يجب عدم الخلط بين:

\[
(fg)(n)=f(n)g(n)
\]
و
\[
(f*g)(n)=\sum_{d\mid n}f(d)g(n/d).
\]

الضرب النقطي يحافظ على الضربية، والالتفاف أيضًا يحافظ عليها، لكنهما عمليتان مختلفتان جذريًا.

\begin{example}
إذا كانت \(1(n)=1\)، فإن
\[
1\cdot1=1،
\]
بينما
\[
1*1=\tau.
\]
\end{example}

\section{الدوال المدعومة على القوى الأولية}

\begin{definition}
نقول إن \(f\) مدعومة على القوى الأولية إذا
\[
f(n)=0
\]
كلما لم يكن \(n=p^k\) لقوة أولية.
\end{definition}

دالة فون مانغولت مثال أساسي:
\[
\Lambda(n)
=
\begin{cases}
\log p,&n=p^k,\\
0,&\text{خلاف ذلك}.
\end{cases}
\]

وهي تحقق الهوية:
\[
\log n
=
\sum_{d\mid n}\Lambda(d).
\]

أي
\[
\log
=
1*\Lambda.
\]

وبقلب موبيوس:
\[
\Lambda
=
\mu*\log.
\]

\section{برهان هوية فون مانغولت}

\begin{theorem}
لكل \(n\geq1\):
\[
\sum_{d\mid n}\Lambda(d)=\log n.
\]
\end{theorem}

\begin{proof}
إذا
\[
n=\prod_{j=1}^{r}p_j^{\alpha_j},
\]
فإن القواسم التي تعطي مساهمة غير صفرية هي القوى
\[
p_j,\ p_j^2,\dots,p_j^{\alpha_j}.
\]
لذلك:
\[
\sum_{d\mid n}\Lambda(d)
=
\sum_{j=1}^{r}\alpha_j\log p_j
=
\log\left(\prod_{j=1}^{r}p_j^{\alpha_j}\right)
=
\log n.
\]
\end{proof}

\section{التفافات مهمة}

من الهويات الأساسية:

\[
1*1=\tau,
\]
\[
1*\operatorname{id}=\sigma,
\]
\[
\mu*1=\varepsilon,
\]
\[
\mu*\operatorname{id}=\varphi,
\]
\[
1*\Lambda=\log.
\]

كما أن:
\[
\tau_k
=
\underbrace{1*1*\cdots*1}_{k\text{ مرات}},
\]
حيث \(\tau_k(n)\) عدد الطرق المرتبة لكتابة \(n\) كحاصل ضرب \(k\) أعداد موجبة.

\section{دوال القواسم العامة}

نعرف:
\[
\sigma_\alpha(n)
=
\sum_{d\mid n}d^\alpha.
\]

إذن:
\[
\sigma_\alpha
=
1*\operatorname{id}_\alpha,
\qquad
\operatorname{id}_\alpha(n)=n^\alpha.
\]

وللقوة الأولية:
\[
\sigma_\alpha(p^k)
=
1+p^\alpha+\cdots+p^{k\alpha}.
\]

إذا \(\alpha\neq0\)، فإن:
\[
\sigma_\alpha(p^k)
=
\frac{p^{(k+1)\alpha}-1}{p^\alpha-1}.
\]

\section{دالة ليوفيل}

\begin{definition}
إذا
\[
n=p_1^{\alpha_1}\cdots p_r^{\alpha_r},
\]
نعرف
\[
\Omega(n)=\alpha_1+\cdots+\alpha_r
\]
و
\[
\lambda(n)=(-1)^{\Omega(n)}.
\]
\end{definition}

الدالة \(\lambda\) ضربية تمامًا، لأن:
\[
\Omega(mn)=\Omega(m)+\Omega(n).
\]

كما ترتبط بموبيوس عبر:
\[
\lambda
=
\mu*1_{\square},
\]
حيث \(1_{\square}(n)=1\) إذا كان \(n\) مربعًا كاملًا، و\(0\) خلاف ذلك.

\section{الدالة \(\omega(n)\) والدالة \(\Omega(n)\)}

نعرف:

\[
\omega(n)=\#\{p:p\mid n\},
\]
أي عدد العوامل الأولية المتميزة.

أما:
\[
\Omega(n)=\sum_{p^\alpha\parallel n}\alpha,
\]
فتحسب العوامل الأولية مع التكرار.

ليستا ضربيتين، لكنهما جمعيتان بمعان مختلفة:

إذا \((m,n)=1\)، فإن
\[
\omega(mn)=\omega(m)+\omega(n).
\]

أما \(\Omega\)، فتتحقق لها الإضافة التامة:
\[
\Omega(mn)=\Omega(m)+\Omega(n)
\]
لكل \(m,n\).

\section{المعكوس الديريشلي للدالة الضربية}

إذا كانت \(f\) ضربية و\(f(1)=1\)، فإن معكوسها يحدد محليًا.

عند قوة أولية:
\[
(f*f^{-1})(p^k)=0
\qquad(k\geq1).
\]

أي:
\[
\sum_{j=0}^{k}
f(p^j)f^{-1}(p^{k-j})=0.
\]

ومنها:
\[
f^{-1}(p^k)
=
-\sum_{j=1}^{k}
f(p^j)f^{-1}(p^{k-j}).
\]

\begin{example}
إذا كانت \(f=1\)، فإن:
\[
1^{-1}=\mu.
\]
\end{example}

\section{المتسلسلات المولدة المحلية}

للدالة الضربية \(f\)، ترتبط معلوماتها المحلية بالمتسلسلة:
\[
\mathcal{B}_p(f;X)
=
\sum_{k=0}^{\infty}f(p^k)X^k.
\]

إذا كان \(f(1)=1\)، فإن:
\[
\mathcal{B}_p(f^{-1};X)
=
\frac1{\mathcal{B}_p(f;X)}.
\]

هذه الهوية المحلية هي النموذج الجبري الذي سيتحول لاحقًا إلى حاصل ضرب أويلري عند وضع:
\[
X=p^{-s}.
\]

\section{التفاف ديريشليه وسلاسل ديريشليه}

إذا كانت:
\[
F(s)=\sum_{n=1}^{\infty}\frac{f(n)}{n^s},
\]
و
\[
G(s)=\sum_{n=1}^{\infty}\frac{g(n)}{n^s},
\]
وكان التقارب مطلقًا، فإن:
\[
F(s)G(s)
=
\sum_{n=1}^{\infty}\frac{(f*g)(n)}{n^s}.
\]

\begin{proof}
بالتقارب المطلق:
\[
F(s)G(s)
=
\sum_{a,b\geq1}
\frac{f(a)g(b)}{(ab)^s}.
\]
نجمع الحدود بحسب \(n=ab\):
\[
F(s)G(s)
=
\sum_{n\geq1}
\frac1{n^s}
\sum_{ab=n}f(a)g(b).
\]
وهذا يساوي:
\[
\sum_{n\geq1}\frac{(f*g)(n)}{n^s}.
\]
\end{proof}

هذه النتيجة تفسر لماذا الالتفاف هو الضرب الطبيعي لمعاملات سلاسل ديريشليه.

\section{البنية الحلقية}

يمكن تلخيص البنية كما يلي:

\[
(\mathcal{A},+,*)
\]
حلقة تبادلية ذات وحدة \(\varepsilon\).

وحداتها هي بالضبط الدوال \(f\) التي تحقق:
\[
f(1)\neq0.
\]

أما الدوال الضربية ذات \(f(1)=1\)، فتشكل زمرة تحت الالتفاف الديريشلي.

\section{ملاحظات تاريخية ومنهجية}

تظهر صيغة قلب موبيوس في القرن التاسع عشر ضمن دراسة علاقات القواسم، لكنها أصبحت لاحقًا نموذجًا عامًا للانقلاب على المرتبات الجزئية وبنى التضمين والاستبعاد.

كما أن الالتفاف الديريشلي لم يعد مجرد أداة حسابية؛ بل صار بنية جبرية مركزية تظهر في:

\begin{itemize}
  \item سلاسل ديريشليه.
  \item دوال زيتا و\(L\).
  \item التوافقيات الجبرية.
  \item نظرية الاحتمالات الحسابية.
  \item التعميمات متعددة المتغيرات.
\end{itemize}

وتوجد تعميمات عديدة للالتفاف، منها الالتفاف الوحدوي، والالتفاف المرتبط بالقاسم المشترك، والالتفافات على أنصاف الزمر الحسابية، لكن الالتفاف الديريشلي سيظل العملية الأساسية في هذه الموسوعة.

\section{أخطاء شائعة}

\begin{enumerate}
  \item نسيان شرط \(f(1)=1\) في تعريف الدالة الضربية.
  \item الخلط بين الضربية والضربية التامة.
  \item الخلط بين الضرب النقطي والالتفاف.
  \item استعمال قلب موبيوس دون التحقق من شكل مجموع القواسم.
  \item افتراض أن كل دالة لها معكوس ديريشلي.
  \item إهمال التقارب المطلق عند ضرب سلاسل ديريشليه.
  \item استنتاج الضربية من قيم قليلة دون إثبات على القوى الأولية.
\end{enumerate}

\section{تمارين}

\begin{enumerate}
  \item أثبت أن حاصل الضرب النقطي لدالتين ضربيتين ضربي.
  \item أثبت أن معكوس الدالة الضربية تحت الالتفاف ضربي.
  \item احسب المعكوس الديريشلي للدالة \(\operatorname{id}\).
  \item أثبت:
  \[
  \sum_{d\mid n}\varphi(d)=n.
  \]
  \item استعمل قلب موبيوس لإثبات:
  \[
  \varphi(n)=n\sum_{d\mid n}\frac{\mu(d)}d.
  \]
  \item أثبت:
  \[
  \sum_{d\mid n}\mu^2(d)=2^{\omega(n)}.
  \]
  \item أثبت أن:
  \[
  \tau_k(p^\alpha)
  =
  \binom{\alpha+k-1}{k-1}.
  \]
  \item أوجد متسلسلة بيل المحلية لـ\(\mu\)، و\(\tau\)، و\(\varphi\).
  \item أثبت:
  \[
  \lambda(n)
  =
  \sum_{d^2\mid n}\mu(n/d^2).
  \]
  \item إذا كانت \(f\) ضربية تمامًا، أثبت أن:
  \[
  f^{-1}(p)=-f(p),
  \]
  وحدد \(f^{-1}(p^2)\) و\(f^{-1}(p^3)\).
  \item أثبت أن دالة فون مانغولت ليست ضربية.
  \item برهن أن:
  \[
  \sum_{d\mid n}\Lambda(d)=\log n
  \]
  بطريقتين مختلفتين.
\end{enumerate}

\section{خلاصة الفصل}

الدوال الحسابية هي اللغة المحلية للبنية الضربية، والالتفاف الديريشلي هو العملية التي تجمع هذه البنية بصورة تتوافق مع ضرب سلاسل ديريشليه. وتؤدي دالة موبيوس دور المعكوس الأساسي، بينما تكشف صيغة قلب موبيوس كيف يمكن استرجاع المعلومات من مجاميع القواسم.

سيبنى الفصل التالي على هذه اللغة لدراسة الدوال الحسابية الأساسية بالتفصيل، قبل الانتقال إلى سلاسل ديريشليه والمنتجات الأويلرية.
